\documentclass[14pt,a4paper]{extarticle}
\usepackage[left=3cm,right=2cm,top=2cm,bottom=2cm,headheight=18pt]{geometry}
\usepackage{fontspec}
\setmainfont{Liberation Serif}
\setsansfont{Liberation Serif}
\setmonofont{Liberation Mono}
\usepackage{unicode-math}
\setmathfont{texgyretermes-math.otf}
\usepackage{amsmath,graphicx,booktabs,longtable,array,enumitem}
\usepackage{titlesec,fancyhdr}
\usepackage[hidelinks,unicode]{hyperref}
\usepackage{url}
\urlstyle{same}
\setlength{\parindent}{1.27cm}
\setlength{\parskip}{0pt}
\linespread{1}
\setlength{\emergencystretch}{2em}
\clubpenalty=10000
\widowpenalty=10000
\titleformat{\section}{\normalfont\bfseries}{\thesection.}{0.5em}{}
\titleformat{\subsection}{\normalfont\bfseries}{\thesubsection.}{0.5em}{}
\titlespacing*{\section}{0pt}{1.2em}{0.5em}
\titlespacing*{\subsection}{0pt}{0.9em}{0.3em}
\setlist{nosep,leftmargin=1.27cm}
\pagestyle{fancy}
\fancyhf{}
\renewcommand{\headrulewidth}{0pt}
\renewcommand{\footrulewidth}{0pt}
\fancyhead[C]{\ifnum\value{page}>1\thepage\fi}
% Liberation Serif is a disclosed Times-compatible substitute, not Times New Roman.
% Current statutory geometry/size baseline: O‘RQ-682 appendix §§43–51.
\hyphenpenalty=10000
\exhyphenpenalty=10000
\pretolerance=10000
\tolerance=10000
\setlength{\emergencystretch}{4em}
\renewcommand{\normalsize}{\fontsize{14bp}{16.8bp}\selectfont}
\normalsize
\begin{document}
\textbf{Yuborilmagan murojaat loyihasi.} 2026-yil 1-oktabrda tahrirlangan. Hisoblash auditining tayanch sanasi 2026-yil 30-sentabr bo‘lib qoladi. Elektron taqdim etish uchun tayyorlangan.

O‘zbekiston Respublikasi Prezidenti huzuridagi Ijtimoiy himoya milliy agentligiga

Nusxa: O‘zbekiston Respublikasi Iqtisodiyot va moliya vazirligiga

Muallif: Abduxoliq Akmal ugli Ashuraliyev

Yashash manzili: Toshkent shahri, Yunusobod tumani, Qashqar mavzesi, 16.

Elektron pochta: \href{mailto:abdulxoliqashuraliyev@gmail.com}{abdulxoliqashuraliyev@gmail.com}.

\section*{PF-193-son Farmon ijrosi bo‘yicha ilmiy-texnik taklif va ishchi guruhda ishtirok etish haqida}
O‘zbekiston Respublikasi Prezidentining 2026-yil 11-sentabrdagi PF-193-son Farmoni 13–14-bandlariga muvofiq ishlab chiqilayotgan normativ-huquqiy hujjatga ko‘p o‘lchovli baholash hisob-kitoblarini tekshirish, ma’lumotlar sifati, hisoblash talqinlarini qayd etish va qaror asoslarini tushuntirish bo‘yicha ilova qilinayotgan mualliflik taklifini taqdim etaman.

Mazmuniy qism “Jismoniy va yuridik shaxslarning murojaatlari to‘g‘risida”gi Qonunning 3-moddasidagi davlat faoliyatini takomillashtirishga doir taklifdir. U PF-193-son Farmon bilan allaqachon topshirilgan loyihaga hissa sifatida taqdim etiladi. “Normativ-huquqiy hujjatlar to‘g‘risida”gi Qonunning 20-moddasi masalani mavjud qonuniy mexanizm va ma’muriy vakolat orqali hal qilish mumkin bo‘lganda loyiha tayyorlashni tashabbus qilishni cheklaydi. Taklif etilgan vazifa mavjud vakolat doirasida qonuniy bajarilsa, tegishli vakolat va amalga oshirish yechimini ko‘rsatishni so‘rayman; amaldagi majburiyatni takrorlash uchungina yangi hujjat so‘ralmaydi. Yangi tashqi majburiyatlar va yuqori yuridik kuchga ega mezonlarning o‘zgarishi vakolatli huquqiy hujjatni talab qiladi.

“Normativ-huquqiy hujjatlar to‘g‘risida”gi Qonunning 20, 22 va 23-moddalariga muvofiq:

\begin{enumerate}[start=1]
\item Taklif etilgan normalarni PF-193-son Farmonning 14-bandi bo‘yicha tayyorlanayotgan loyiha doirasida, avvalo besh bandli qisqa muqobil yoki teng normativ yechimni ko‘rib chiqishingizni;
\item Mazkur yo‘nalishga mas’ul loyiha tayyorlovchi bo‘linma yoki ishchi guruh bilan ilmiy-texnik muhokama tashkil etish masalasini ko‘rib chiqishingizni;
\item Quyidagi bajarilgan ishlar va ilovadagi takrorlash vositalari asosida meni hisoblash uslubiyotini tekshirish yo‘nalishida ishchi guruh tarkibiga kiritish imkoniyatini ko‘rib chiqishingizni so‘rayman. Ushbu ishtirokka rozilik bildiraman. Tegishli bilim va tajribani baholash uchun texnik muhokama hamda hisob-kitob namoyishida ishtirok etishga tayyorman;
\item Joriy loyiha va sanali spetsifikatsiyani tasdiqlovchi akt yoki loyiha maqomi, birlik, davr, chegara, istisno, yaxlitlash va sifat qoidalari bilan taqdim etish imkoniyatini ko‘rib chiqishingizni. Qonuniy ochiq qismda takroriy ishni oldini olish uchun mavjud mustaqil kutilgan javob sinovi, chiqarilish/kelishuv qaydi va ilova misoli qamrovini ko‘rsatishni so‘rayman. Qoplanmagan vazifa uchun mas’ul, shartnoma qamrovi, qo‘shimcha ish/soat, qiymat tartibi va qonuniy mablag‘ belgilansin. Ariza ishlovini o‘zgartirishda haqiqiy yuk va vaqt ishlatiladi; ochiq hisoblagich quvvat bahosi sifatida taklif etilmaydi. Oila yozuvi va himoyalangan xavfsizlik tafsiloti so‘ralmaydi;
\item Har bir taklif bo‘yicha loyiha matniga kiritilgan tahrirni, o‘zgartirib qabul qilingan yechimni, masala amaldagi normada hal qilingan bo‘lsa uning aniq asosini yoki taklif inobatga olinmagan bo‘lsa sababini ko‘rsatishingizni so‘rayman. Boshqa hujjatda hal qilinadigan masala uchun tegishli hujjat va mas’ul organni ko‘rsatish foydali bo‘ladi. Bu band fuqaro taklifiga so‘ralayotgan javob shaklidir; 24-modda bo‘yicha portal muhokamasi o‘tkazilgan deb ko‘rsatilmaydi.
\end{enumerate}
22-modda doirasida taklif etayotgan hissam tasdiqlangan uslubiyotni takror hisoblash mumkin bo‘lgan qoidalar tavsifiga o‘tkazish, chegaralar va istisnolar uchun nazorat holatlarini tayyorlash, birlik va davrlarni tekshirish hamda mas’ul xodimlar bilan tegishli tekshiruv normalarini ishlab chiqish bilan cheklanadi. Murojaatni PF-193-son Farmon bo‘yicha loyihani tayyorlayotgan bo‘linmaga, zarur hollarda Agentlik huzuridagi “Axborot texnologiyalar markazi” davlat muassasasining texnik ishtirokini ta’minlagan holda yo‘naltirishni so‘rayman. 35-son qarorning 1-ilovasi 46-bandida Reyestr va axborot tizimining qonunchilik hamda texnik hujjatlarga muvofiq to‘g‘ri ishlashi uchun shu muassasa mas’ullari javobgar etib belgilangan. Yordam siyosatini tanlash yoki huquqiy mezonlarni tasdiqlash vakolatini so‘ramayman. Ishchi guruhga kiritish tanlanmasa, 22-modda doirasida ilovadagi ilmiy-texnik tavsiyalarni loyiha tayyorlashda qo‘llash imkoniyatini ko‘rib chiqish va keyingi namoyish uchun mas’ul aloqa shaxsini ko‘rsatishni so‘rayman.

Asosiy tekshiruv ilovasi tushuntirish xatidagi toifalash tengligi va yer formulasi bo‘yicha ixtiyoriy o‘zgartirish variantlaridan alohida ko‘rib chiqilishi mumkin. Bu variantlar vakolatli huquqiy o‘zgartirish hamda oqibatlarni baholashni talab qiladi; ular tekshiruv taklifini qabul qilish sharti emas.

1–2-material zarur normativ mazmunni ichki tekshiruvdan ajratadi. Avvalo 1-materialning besh qisqa muqobil bandini PF-193 asosiy loyihasida yoki vakolat, soha va sanani saqlaydigan teng normada ko‘rishni so‘rayman. To‘liq tahrir birlashtirish varianti, qo‘shimcha jamlanadigan talab emas. Texnik vazifaga teng mavjud sinov ishlatilsin; huquqiy vazifaga shu majburiy oqibatni ta’minlaydigan aniq norma ko‘rsatilsin. Sinov qamrovi yo‘q huquqiy mezonni belgilamaydi. Vakolatli kirish cheklovi yoki mavjud qoidaning sharhi manba, soha va sana bilan qabul qilinadi; sharh yangi mezon kiritmaydi. Qisqa normativ matn yoki havola mazmuniy qabul natijasidir. Respublika qo‘lda xizmati, yangi statistik hisobot, 2026 uslubiyotini saqlash, o‘tish huquqi va qarzdan ozod qilish taklif etilmaydi. Qo‘shimcha soat, shartnoma qamrovi, quvvat va mablag‘ tasdiqlanmagan; faqat qoplanmagan integratsiya baholanadi.

Murojaatni “Jismoniy va yuridik shaxslarning murojaatlari to‘g‘risida”gi Qonunning 23-moddasi bo‘yicha ro‘yxatdan o‘tkazib, kelib tushgan sana va ro‘yxat raqamini ma’lum qilishni so‘rayman. Ushbu modda ro‘yxatdan o‘tkazishni rad etishni taqiqlaydi; 18-modda 29–30-moddalardagi holatlardan tashqari qabul qilish va ko‘rib chiqishni talab qiladi. Masala oluvchining vakolatiga kirmasa, murojaatlar to‘g‘risidagi qonunga muvofiq besh kunlik muddatda vakolatli organga yuborib, bu haqda yozma yoki elektron xabar berishni so‘rayman. Bir necha ichki bo‘linmaga tegishli masalalarni 24-modda bo‘yicha muvofiqlashtirishni so‘rayman. Bu havolalar \href{https://lex.uz/docs/-3336169}{murojaatlar haqidagi Qonunga} tegishli, norma ijodkorligi haqidagi Qonunning boshqa mazmundagi moddalari bilan tenglashtirilmaydi.

Murojaatlar haqidagi Qonunning 27-moddasi bo‘yicha har bir masalaga aniq asoslar bilan yozma yoki elektron javob yuborishni so‘rayman. 28-moddada normativ taklif kelib tushgan kundan e’tiboran bir oygacha ko‘rib chiqilishi, qo‘shimcha o‘rganish talab etilsa taklif kiritgan shaxsga o‘n kunlik muddatda yozma xabar berilishi belgilangan. Shaxsiy ishtirok va kirish haqidagi iltimoslar alohida ko‘rib chiqilsa, ularga tatbiq etilgan murojaat turi va muddatni ko‘rsatishni so‘rayman; boshqa turdagi murojaatga taklif muddatini majburan tatbiq etmayman. Taklif yoki iltimos qabul qilinmasa, 33-modda bo‘yicha tegishli shikoyat qilish tartibini tushuntirishni so‘rayman. Ishchi guruhga kiritish yoki hujjatni qabul qilishga talab huquqi da’vo qilinmaydi.

\href{https://lex.uz/docs/-5378966}{“Normativ-huquqiy hujjatlar to‘g‘risida”gi Qonun}ning 23-moddasi ishlab chiquvchidan tegishli ilmiy tadqiqot va fuqarolar murojaatlarini o‘rganish, taklif hamda ilmiy tavsiyalarni umumlashtirish va loyiha tayyorlashda foydalanishni talab qiladi. Ishchi guruh ishtiroki tanlanmasa ham ilovadagi natijalarni shu jarayonda ko‘rib chiqishni so‘rayman. Jamoatchilik yoki mutaxassislar muhokamasida qatnashsam, qonuniy cheklovlarga rioya etilgan holda ushbu modda bo‘yicha loyiha matni bilan oldindan tanishtirishni so‘rayman. Bu himoyalangan axborotga cheklanmagan kirish huquqi yoki portal muhokamasi allaqachon o‘tkazilganligi haqidagi da’vo emas.

Haqiqiy oilalarning shaxsga doir ma’lumotlarini menga yuborish so‘ralmaydi. Hisoblash vositalarini vakolatli operatorning nazoratidagi muhitda, qonuniy asos va egasizlantirish talablariga rioya etgan holda qo‘llash taklif qilinadi. Sintetik misollar dasturiy tekshiruvni namoyish etadi; ular haqiqiy oilalar toifalari yoki ehtiyojini aniqlash to‘g‘riligi yuzasidan dalil emas. Amaldagi ochiq normalar bo‘yicha hisoblash auditi 35-son qarordagi rasmiy misolni takrorlaydi, daromad aynan minimal iste’mol xarajatiga teng holatdagi toifalash shartlari va manfiy hosila yer maydoni uchun nazorat holatlarini ko‘rsatadi. Bu holatlarning amaldagi axborot tizimida uchrashi yoki ularga ta’sirlangan oilalar soni hali aniqlanmagan.

22-modda bo‘yicha tayyorlashda ishtirok etishga roziman. Taklif etilgan hissa — normativ matn va qayta bajariladigan nazorat hisob-kitoblari.

Ilovalar mazmuniy ko‘rib chiqish tartibida:
\begin{enumerate}
\item 1-material — «Oilalarni ko‘p o‘lchovli baholash hisob-kitoblarini tekshirishga doir normativ qoidalar loyihasi»: mavjud qonuniy nazoratga integratsiya uchun yetti asosiy band va 28-bandli tartib. Asosiy ruscha tahrir va mos o‘zbekcha hamda inglizcha tarjimalar. To‘liq shaklni almashtiradigan besh bandli qisqa muqobil ham ilova qilingan.
\item 2-material — «Tushuntirish xati»: topshiriq va normalar zarurati, bandlar bo‘yicha ijro asosi, hisoblangan natijalar, tanlangan qiyosiy jadval, javobgarlik, resurs hisobi va oqibatlar. Asosiy ruscha tahrir va mos o‘zbekcha hamda inglizcha tarjimalar.
\item 3-material — «Ilmiy va huquqiy asoslantirish»: usul, natija va dalillar, manbalar va xulosa chegaralari. Elektron hisoblash qo‘shimchasi oldindan belgilangan reja, sanali qoida tavsifi, nazorat holatlari, dastur manbalari, mustaqil arifmetik tekshiruv va bajarish natijalarini qamraydi. Asosiy ruscha tahrir va mos o‘zbekcha hamda inglizcha tarjimalar.
\end{enumerate}
Uchta material va hisoblash qo‘shimchasi murojaat bilan birga taqdim etiladi. Mazmun va hisoblash xulosalarini repozitoriyga kirmasdan ko‘rib chiqish mumkin. Asosiy murojaat tili — ruscha; murojaatlar to‘g‘risidagi qonunning 6-moddasi bunga ruxsat beradi. O‘zbekcha tahrir O‘RQ-682-son Qonunning 21-moddasi bo‘yicha rasmiy loyihani davlat tilida tayyorlash uchun beriladi. Elektron taqdim etishga sanalgan materiallarning joriy tahrirlari ilova qilinadi.
\section*{22-modda bo‘yicha tegishli ishlar bayoni}
Taqdim etilayotgan loyiha doirasida amaldagi ochiq daromad formulalarining aniq ratsional arifmetika bilan hisoblash vositasi, alohida hisoblash tekshiruvi va avtomatlashtirilgan sinovlar tayyorlangan. Tekshiruv 37 ta sinovdan iborat; yer formulasi uchun 4\,105 ta tuzilgan holat va 620 ta manfiy hosila maydon holati qayd etilgan. Rasmiy misol qayta hisoblanganda 432\,000 so‘m natija olingan. Bu raqamlar haqiqiy oilalarga doir tanlanma natijalari emas.

Elektron hisoblash qo‘shimchasi bajarish tartibi, aniq kirish qiymatlari, kutilgan javoblar, misollar tarkibining mustaqil tekshiruvi va ko‘rsatilgan raqamlar asosidagi natijalarni qamraydi. U 22-modda uchun tegishli bajarilgan ishni tasdiqlaydi; ilmiy daraja yoki tashkilot vakilligi da’vo qilinmaydi. Hisoblarni loyiha tayyorlash jamoasiga namoyish etishga tayyorman.

Abduxoliq Akmal ugli Ashuraliyev\par
Tahrir sanasi: 2026-yil 1-oktabr.

\end{document}
